\documentclass[11pt,a4paper]{article}
\usepackage[margin=2.5cm]{geometry}
\usepackage{amsmath,amssymb}
\usepackage{booktabs}
\usepackage{hyperref}
\usepackage{listings}
\usepackage{xcolor}
\usepackage{parskip}

\lstset{
  basicstyle=\ttfamily\small,
  breaklines=true,
  frame=single,
  backgroundcolor=\color{gray!10}
}

\title{\textbf{Introduction to Data Mining, Spring 2026}\\[4pt]
       \Large HW1 Write-up}
\author{}
\date{}

\begin{document}
\maketitle
\vspace{-3em}
\noindent Student ID: 202419212
\bigskip

% ============================================================ %
\section*{Use of AI Tools}
% ============================================================ %

Claude Code (Anthropic, \texttt{claude-sonnet-4-6}) was used in completing this homework.
Specifically, it assisted with (1) writing and debugging all four Python programs
(\texttt{hw1\_1.py}, \texttt{hw1\_2\_b.py}, \texttt{hw1\_3\_b.py}, \texttt{hw1\_3\_c.py}),
and (2) drafting this write-up.
The algorithmic ideas, design-choice justifications, and mathematical derivations for
the textbook exercises were developed jointly through interactive dialogue.
All code was reviewed and verified by the student.

% ============================================================ %
\section{Problem 1 — MapReduce and Spark}
% ============================================================ %

\subsection*{Algorithm Description}

The goal is to find pairs of users who share many mutual friends but are not
direct friends.
The algorithm has five logical steps, each implemented as a Spark RDD transformation.

\begin{enumerate}
  \item \textbf{Parse.} Each line is split on the tab character to produce
        \texttt{(user:\ int,\ friends:\ list[int])}.

  \item \textbf{Emit mutual-friend candidates.}
        For each user $u$ with friend list $F$, every unordered pair
        $\{f_i, f_j\} \subseteq F$ has $u$ as a mutual friend.
        We emit $\bigl((\min(f_i,f_j),\,\max(f_i,f_j)),\;1\bigr)$
        for all such pairs.
        Canonicalising by $({\rm min},{\rm max})$ ensures each pair is counted
        once regardless of which endpoint emits it.

  \item \textbf{Emit existing friendships.}
        For each user $u$ and each of their friends $f$, emit
        $\bigl((\min(u,f),\,\max(u,f)),\;{\rm True}\bigr)$.
        These keys mark direct-friend pairs that must be excluded.

  \item \textbf{Count and filter.}
        \texttt{reduceByKey} sums the mutual-friend counts.
        \texttt{subtractByKey} removes any pair that appears in the existing-friendship
        RDD.

        \textit{Design choice — \texttt{subtractByKey} vs.\ broadcast filter:}
        The full friendship set can be tens of millions of pairs and would not fit
        comfortably in a broadcast variable on every executor.
        \texttt{subtractByKey} performs a shuffle-based join, which Spark handles
        efficiently without requiring the set to reside in driver memory.

  \item \textbf{Sort and output.}
        Results are sorted by count (descending), then first user ID (ascending),
        then second user ID (ascending), and the top~10 are taken.
\end{enumerate}

\textit{Design choice — in-Python nested loop for pair generation:}
Because \texttt{itertools} is not a permitted library, pairs are generated with a
double \texttt{for} loop.
This is $O(|F|^2)$ per user, which is acceptable because the LiveJournal adjacency
list has moderate degree for most users.

\subsection*{Top-10 Output}

\begin{lstlisting}
18739   18740   100
31506   31530   99
31492   31511   96
31511   31556   96
31519   31554   96
31519   31568   96
31533   31559   96
31555   31560   96
31492   31556   95
31503   31537   95
\end{lstlisting}

\subsection*{Elapsed Time}

Measured on a local WSL2 machine (Intel Core, PySpark 3.5.1, Java 1.8):
\textbf{168 seconds} ($\approx 2.8$ minutes).

% ============================================================ %
\section{Problem 2 — Frequent Itemsets}
% ============================================================ %

\subsection*{(a) Exercise 6.2.7}

\textbf{Pass~1} counts all $10^6$ items ($4\,\text{MB}$).

\textbf{Pass~2 — Option~A (triangular matrix):}
$\binom{N}{2}$ candidate pairs, 4 bytes each:
\[
  M_A = 2N(N-1) \text{ bytes.}
\]

\textbf{Pass~2 — Option~B (hash table of triples):}
$10^6 + M$ non-zero pairs of frequent items, 12 bytes each:
\[
  M_B = 12(10^6 + M) \text{ bytes.}
\]

\textbf{Minimum memory} (Pass~2 dominates):
\[
  \boxed{\min\!\bigl(2N(N-1),\; 12(10^6 + M)\bigr) \text{ bytes.}}
\]

\subsection*{(b) A-Priori Implementation}

\begin{itemize}
  \item \textit{Pass~1.} Count items with a \texttt{dict} ($O(1)$ update, no pre-allocation needed).

  \item \textit{Pass~2 (triangular matrix).} Project each basket onto frequent items,
        deduplicate, and increment the \texttt{numpy.int32} triangular array for every pair.
        Index formula for $i < j$:
        \[
          \texttt{index}(i,j) = i N - \tfrac{i(i+1)}{2} + (j - i - 1).
        \]

  \item \textit{Association rules.} For each frequent pair $(A,B)$, generate $A\!\to\!B$
        and $B\!\to\!A$; keep confidence $\ge 0.5$; sort by confidence then support (descending).
\end{itemize}

\textbf{Output (12 lines):}
\begin{lstlisting}
1334
39
Rule: DAI93865 -> FRO40251, Confidence: 1.000000, Support: 208
Rule: GRO85051 -> FRO40251, Confidence: 0.999176, Support: 1213
Rule: GRO38636 -> FRO40251, Confidence: 0.990654, Support: 106
Rule: ELE12951 -> FRO40251, Confidence: 0.990566, Support: 105
Rule: DAI88079 -> FRO40251, Confidence: 0.986726, Support: 446
Rule: FRO92469 -> FRO40251, Confidence: 0.983510, Support: 835
Rule: DAI43868 -> SNA82528, Confidence: 0.972973, Support: 288
Rule: DAI23334 -> DAI62779, Confidence: 0.954545, Support: 273
Rule: ELE92920 -> DAI62779, Confidence: 0.732665, Support: 877
Rule: DAI53152 -> FRO40251, Confidence: 0.717949, Support: 140
\end{lstlisting}

\textbf{Elapsed time:} \textbf{0.73 seconds.}

% ============================================================ %
\section{Problem 3 — Finding Similar Items}
% ============================================================ %

\subsection*{(a) Exercise 3.6.1}

Let $s \in [0,1]$ be the Jaccard similarity (probability that a single hash function agrees).
An $r$-way AND of inputs with probability $p$ yields $p^r$;
a $b$-way OR yields $1-(1-p)^b$.

\begin{enumerate}
  \item[(a)] \textbf{2-way AND $\to$ 3-way OR.}
        After AND: $p_1 = s^2$.
        After OR:
        \[ p = 1 - (1 - s^2)^3. \]

  \item[(b)] \textbf{3-way OR $\to$ 2-way AND.}
        After OR: $p_1 = 1-(1-s)^3$.
        After AND:
        \[ p = \bigl(1-(1-s)^3\bigr)^2. \]

  \item[(c)] \textbf{2-way AND $\to$ 2-way OR $\to$ 2-way AND.}
        \begin{align*}
          p_1 &= s^2, \\
          p_2 &= 1-(1-s^2)^2, \\
          p   &= \bigl(1-(1-s^2)^2\bigr)^2.
        \end{align*}

  \item[(d)] \textbf{2-way OR $\to$ 2-way AND $\to$ 2-way OR $\to$ 2-way AND.}
        \begin{align*}
          p_1 &= 1-(1-s)^2, \\
          p_2 &= \bigl(1-(1-s)^2\bigr)^2, \\
          p_3 &= 1-\bigl(1-p_2\bigr)^2
               = 1-\Bigl(1-\bigl(1-(1-s)^2\bigr)^2\Bigr)^2, \\
          p   &= p_3^2
               = \Bigl(1-\Bigl(1-\bigl(1-(1-s)^2\bigr)^2\Bigr)^2\Bigr)^2.
        \end{align*}
\end{enumerate}

\subsection*{(b) MinHash-Based LSH}

\begin{itemize}
  \item Shingle extraction: lowercase letters and spaces only; 3-character sliding window.
        Each article's shingles are stored in a Python \texttt{set} for $O(1)$ deduplication and lookup.

  \item Global shingle list sorted lexicographically to ensure a deterministic row-index mapping.

  \item MinHash ($b=6$, $r=20$, 120 hash functions): $h_i(x)=(a_i x+b_i)\bmod c$,
        $c$ = smallest prime $\ge n$, $a_i,b_i \sim \mathrm{Uniform}[0,c-1]$.
        All 120 hash values per document are computed as one $(120\times|\text{shingles}|)$
        matrix operation (vectorised with numpy), avoiding a Python loop over hash functions.

  \item LSH: for each band, documents are bucketed by a tuple of 20 signature values
        (hashable dict key). A \texttt{set} accumulates candidate pairs across all bands,
        eliminating duplicates automatically.

  \item Output: pairs with signature similarity $= (\text{matching positions})/120 \ge 0.9$.
\end{itemize}

\textbf{Output:}
\begin{lstlisting}
t8413   t269    1.000000
t980    t2023   0.991667
t3268   t7998   0.991667
t1621   t7958   1.000000
t448    t8535   1.000000
\end{lstlisting}

\textbf{Elapsed time:} \textbf{1.68 seconds.}

\subsection*{(c) TF-IDF and Cosine Distance}

\begin{itemize}
  \item Tokenisation: lowercase and split on whitespace (no punctuation stripping,
        to avoid ambiguous word boundaries).

  \item TF-IDF: $\mathrm{TF}(t,d)=$ raw count; $\mathrm{IDF}(t)=\ln(N/\mathrm{df}(t))$.
        Each document is stored as a sparse \texttt{dict\{term\_idx: tfidf\}} rather than
        a dense vector ($D\!\times\!V \approx 400$\,MB float64 otherwise).
        Global vocabulary sorted lexicographically for a deterministic mapping.

  \item LSH (random hyperplanes): 10 planes from $\mathcal{N}(0,1)$ with \texttt{np.random.seed(0)},
        hash bit $k$ = sign$(\tilde{v}_d \cdot h_k)$, giving a 10-bit bucket key.
        Projections are accumulated term-by-term to avoid materialising the dense $D\!\times\!V$ matrix.

  \item Cosine similarity for candidate pairs in the same bucket:
        \[
          \cos(a,b) = \frac{\tilde{v}_a \cdot \tilde{v}_b}{\|\tilde{v}_a\|\,\|\tilde{v}_b\|}.
        \]
        Dot product iterates over the shorter dict — $O(\min(|d_a|,|d_b|))$ instead of $O(V)$.
        Pairs with cosine distance $1-\cos(a,b) < 0.1$ are output.
\end{itemize}

\textbf{Output:}
\begin{lstlisting}
t8413   t269    1.000000
t980    t2023   0.997668
t1621   t7958   1.000000
t448    t8535   1.000000
\end{lstlisting}

\textbf{Elapsed time:} \textbf{6.45 seconds.}

\end{document}
